% 数理物理1 数理物理1-12 問題1（092-bunsan。当時は 数理物理1-10 問題1・shou10-l1-bunsan） 解いた日 2026-10-21
%% kotae
\toi{1}{○}{立式}
\begin{align*}
  \omega_\mu&=\sqrt{\frac{2K}{m}\cdot2\sin^{2}\Bigl(\frac{\pi\mu}{N}\Bigr)}\\
  &=2\sqrt{\frac{K}{m}}\,\ayamari{\sin\Bigl(\frac{\pi\mu}{N}\Bigr)}{\abs*{\sin\Bigl(\frac{\pi\mu}{N}\Bigr)}}\\
  &=2\sqrt{\frac{K}{m}}\sin\Bigl(\frac{k_\mu a}{2}\Bigr)
\end{align*}
\naze{$\omega_\mu\ge0$ なので，根号を外すと $\abs{\sin}$ になる。$k_\mu a>\pi$（$\mu>N/2$）では $\sin(k_\mu a/2)<0$ である。正しくは $\omega_\mu=2\sqrt{K/m}\,\abs{\sin(k_\mu a/2)}$。以降はこの誤りを引き継ぐ。}
\toi{2}{◎}{}
\begin{align*}
  \sin\Bigl(\frac{k_\mu a}{2}\Bigr)&=\frac{\cos0}{1!}\Bigl(\frac{k_\mu a}{2}\Bigr)-\frac{\cos0}{3!}\Bigl(\frac{k_\mu a}{2}\Bigr)^{3}\\
  &\qquad+\frac{\cos0}{5!}\Bigl(\frac{k_\mu a}{2}\Bigr)^{5}\\
  &\simeq\frac{k_\mu a}{2}
\end{align*}
\[ \omega_\mu=ak_\mu\sqrt{\frac{K}{m}},\qquad v_s=a\sqrt{\frac{K}{m}} \]
\toi{3}{○}{単位}
\[ v_s=4564\ldots=\qty{4600}{\metre\per\second} \]
\[ \sin\Bigl(\frac{k_\mu a}{2}\Bigr)=1\qquad\therefore\ k_\mu a=\pi \]
\begin{align*}
  \omega_{\max}&=2\cdot1\cdot\sqrt{\frac{K}{m}}=3.65\ldots\times10^{13}\\
  &=3.7\times10^{13}\ \nuke{\unit{\per\second}}\\
  \nu_{\max}&=\frac{\omega_{\max}}{2\pi}=5.8\times10^{12}\ \nuke{\unit{\hertz}}
\end{align*}
\naze{数値は正しいが，$\omega_{\max}$ に \unit{\per\second}，$\nu_{\max}$ に \unit{\hertz} の単位が無い。}
\toi{4}{△}{計算}
\begin{align*}
  \frac{\omega_\mu}{v_sk_\mu}&=\frac{2\sin(1/2)\sqrt{K/m}}{a\sqrt{K/m}\cdot\frac{1}{a}}=2\sin\Bigl(\frac{1}{2}\Bigr)\\
  &=0.9588\ldots
\end{align*}
$k_\mu a=t$ とする。
\begin{align*}
  \frac{\omega_\mu}{v_sk_\mu}&=\frac{2\sin(t/2)\sqrt{K/m}}{a\sqrt{K/m}\cdot\frac{t}{a}}=\frac{2}{t}\sin\Bigl(\frac{t}{2}\Bigr)\\
  &\simeq\ayamari{2\cdot\frac{t}{2}}{\frac{2}{t}\cdot\frac{t}{2}}-\frac{1}{3}\cdot\frac{t^{2}}{8}\\
  &=t-\frac{1}{24}t^{2}=0.99
\end{align*}
\[ \therefore\ t^{2}-24t+\frac{594}{25}=0 \]
\[ \therefore\ t=1.0345\ldots\approx1.03 \]
\begin{align*}
  \lambda_\mu&=v_s\frac{2\pi}{\omega_\mu}=\frac{a\sqrt{K/m}\cdot2\pi}{2\sin(t/2)\sqrt{K/m}}\\
  &=\frac{\pi a}{\sin(t/2)}
\end{align*}
6.35 倍
\naze{$\frac{2}{t}\bigl(\frac{t}{2}-\frac{t^{3}}{48}\bigr)$ の第1項は $\frac{2}{t}\cdot\frac{t}{2}=1$ で，$t$ ではない（第2項には $\frac{2}{t}$ をかけて $t^{2}/24$ にしている）。正しくは $1-t^{2}/24=0.99$ から $t^{2}=0.24$，$k_\mu a\approx0.49$，$\lambda_\mu=2\pi/k_\mu=2\pi a/t\approx13a$。以降はこの誤りを引き継ぐ。また $\lambda_\mu=v_s\cdot2\pi/\omega_\mu$ は，$\omega_\mu/k_\mu$ が $v_s$ と 1\% ずれるので近似式である。定義どおり $\lambda_\mu=2\pi/k_\mu$ を使えば，求めた $t$ からすぐ出る。}
%% kizuki
\item \mae{数理物理1-10（組み替え前）問題1 (1)・10-07} 前回も根号を外すときに絶対値を落とした（そのときは減点なし）。10-03 の答案では $\abs{\sin}$ と書けていた。$\sqrt{\sin^{2}\theta}=\abs{\sin\theta}$ で，$\omega_\mu\ge0$ だから絶対値が要る。根号を外したら，$\mu$ の範囲（$k_\mu a$ は $0$〜$2\pi$）で符号が変わらないかを見る。
\item \mae{数理物理1-10（組み替え前）問題1 (3)・10-07} 前回と同じく $\omega_{\max}$ と $\nu_{\max}$ の単位が無い（今回は $v_s$ だけ付けた）。数値の答えを書いたら，3つとも単位（\unit{\metre\per\second}，\unit{\per\second}，\unit{\hertz}）が付いているかを最後に見る。
\item \mae{数理物理1-10（組み替え前）問題1 (4)・10-07} 前回は比の式で分母の $x$ を落とし，今回は分母をかけたときに第1項だけ $t$ が残った。式を立てたら $t=1$ を入れて前半の値 $0.96$ に戻るかを確かめる（$t-t^{2}/24$ に $t=1$ を入れると $0.958$ で一見合うので，$t$ を小さくして比が 1 に近づくかも見る。$t\to0$ で $t-t^{2}/24\to0$ になるのはおかしい）。$k_\mu a=1$ で $0.96$ なのに $0.99$ になる $k_\mu a$ が 1 より大きいのも食い違いの合図である。
